\section{Extremal characters and invariant tensors}
\label{ext:section}

The tensor conversion proved above is ordinary: its input tensors need
not be invariant.  This permits another way to handle compact
monodromy.  We first construct invariant tensors on the compact factor,
then descend them before applying Corollary~\ref{conv:tensor-conversion}.
The character construction below uses the extreme terms of the full
equivariant Euler characteristic.  It is independent of the
weight-zero index calculation used in the preceding applications.

\begin{proposition}\label{ext:invariant-tensors}
Let $F$ be a smooth connected projective complex variety such that
$H^0(F,\Omega_F^i)=0$ for every $i>0$.  Suppose $L_F=-K_F$ has a
smooth semipositive Hermitian metric.  Let an algebraic torus $T$ act
effectively on $F$, and give $L_F$ its natural tangent-determinant
linearization.  Then there are an integer $p\geq0$ and unbounded
positive integers $k$ such that
\[
 H^0\bigl(F,(\Omega_F^1)^{\otimes p}\otimes kL_F\bigr)^T\ne0.
\]
The trivial torus and the point variety are included.
\end{proposition}

We use the left action on sections: a group element transports a
fiber value and changes its argument by the inverse action.  Thus a
tangent weight $a$ induces conormal weight $-a$.  This convention
fixes both the orbit-degree identity and the expansions below.  The
use of torus characters follows the equivariant viewpoint in
\cite{Muller2025}; the present argument produces tensors to be
converted on the total space, without a semiampleness assumption on
the compact factor.

\begin{proof}
Assume first that $T$ is nontrivial.  Choose a general lattice
cocharacter $u:\mathbb C^*\to T$ for which $F^u=F^T$.  A
$T$-linearized projective embedding shows that such choices exist:
one avoids the finitely many hyperplanes given by differences of its
weights.  The fixed locus has finitely many smooth connected
components $P$.  Write $w_P\in X^*(T)$ for the character of
$L_F|_P$, and put $\lambda_P=\langle w_P,u\rangle$.

\emph{The source has maximal anticanonical weight.}
The Bia\l ynicki-Birula decomposition~\cite{BB1973} has a unique
component $P_s$ whose attracting set as $t\to0$ is open.  The
attracting and repelling sets are affine-space bundles with ranks
given by the positive and negative normal weights; see also
\cite[Theorem~1.5 and Corollary~7.3]{JelisiejewSienkiewicz2019}.
Every normal weight at $P_s$ is positive on $u$.  The component is
proper because the action is nontrivial and effective.  The natural
anticanonical character is the sum of its normal tangent characters,
so
\begin{equation}\label{ext:source-positive}
 \lambda_{P_s}>0.
\end{equation}
For every other component there is a negative normal weight.  Its
repelling cell supplies a nonconstant orbit with limit in that
component as $t\to\infty$.  Extend the \emph{parametrized} orbit to
$a:\mathbb P^1\to F$ by properness.  For a linearized line $B$ the
fiber weights at its two endpoints satisfy
\begin{equation}\label{ext:orbit-degree}
 \operatorname{wt}_u(B|_{a(0)})-
 \operatorname{wt}_u(B|_{a(\infty)})
       =\deg a^*B.
\end{equation}
For example this follows by comparing equivariant frames at zero
and infinity on $\mathbb P^1$.  The parametrization, rather than
only the normalization of its image, includes any finite stabilizer
multiplicity.  For an ample linearized $B$ the difference is
positive, and for the nef line $L_F$ it is nonnegative.  Iterating
backward through fixed components terminates, because the ample
weights strictly increase in a finite set.  It can terminate only
at the source.  Therefore
\begin{equation}\label{ext:source-maximum}
 \lambda_{P_s}\geq\lambda_P\quad\hbox{for all }P.
\end{equation}

The attracting-cell limit map is a dominant rational map
$F\dashrightarrow P_s$.  Resolve it.  Pullback of holomorphic
forms along a dominant map is injective in characteristic zero,
and holomorphic forms are unchanged under smooth birational
modification.  Consequently $P_s$ has no positive-degree
holomorphic forms.  Since $P_s$ is connected and projective,
Hodge symmetry gives
\begin{equation}\label{ext:source-euler}
 \chi(P_s,\mathcal O_{P_s})=1.
\end{equation}

\emph{Localization detects the complete source character.}
Apply fixed-component localization to the virtual character
\[
 \chi_T(F,\ell L_F)=
       \sum_q(-1)^q[H^q(F,\ell L_F)]
\]
for positive integers $\ell$; we use the holomorphic index formula
of \cite{AtiyahSegal1968,AtiyahSinger1968}.  On a fixed component,
write a tangent normal block as $N_a e^a$, where $N_a$ is its
underlying rank-$r$ bundle with the weight removed.  Its contribution
is the factor $\lambda_{-1}(N_a^*e^{-a})^{-1}$.  Expand in the
direction of decreasing pairing with $u$.  If
$\langle a,u\rangle>0$, this factor is
\[
 \sum_{j\geq0}\operatorname{Sym}^j(N_a^*)e^{-ja};
\]
the leading term is $1$ and all others have negative pairing.
If $\langle a,u\rangle<0$, the same factor is
\[
 (-1)^r\det(N_a)e^{ra}
       \sum_{j\geq0}\operatorname{Sym}^j(N_a)e^{ja};
\]
even its leading term has strictly negative pairing.  These
identities follow by factoring the highest exterior term in the
second case.  Every normal character pairs nontrivially with $u$,
so each fixed exponent receives only finitely many contributions;
equivalently one works in the descending Laurent completion over
the character field of the sublattice pairing to zero.

By \eqref{ext:source-maximum}, every nonsource contribution is
strictly below $\ell\lambda_{P_s}$, including components with
$\lambda_P=\lambda_{P_s}$.  At the source the only term at that
pairing is the constant term of each positive normal factor.
Thus the entire part at the highest pairing is
\begin{equation}\label{ext:highest-character}
 e^{\ell w_{P_s}}\,
       \chi(P_s,\ell L_F|_{P_s}).
\end{equation}
This records the full $T$-character, not merely its restriction
to the chosen one-parameter subgroup.  Riemann--Roch makes its
coefficient a polynomial in $\ell$ whose value at zero is $1$
by \eqref{ext:source-euler}.  The coefficient is therefore
nonzero for all sufficiently large positive $\ell$.  Hence the
character $\ell w_{P_s}$ occurs in at least one actual cohomology
group $H^q(F,\ell L_F)$ for each such $\ell$.

\emph{Canceling source characters produces tensors.}
Only finitely many fixed components can occur as sources while
$u$ varies among general lattice directions.  Zero lies in the
convex hull of their characters.  Otherwise strict rational
separation, followed by a small general perturbation, gives a
lattice direction pairing negatively with every source character.
Its own source contradicts \eqref{ext:source-positive}.  Since
the characters are integral, there are nonnegative integers $b_s$,
not all zero, with
\begin{equation}\label{ext:cancel}
 \sum_s b_sw_{P_s}=0.
\end{equation}
For each source with $b_s>0$, choose a general direction having
that source.  Formula~\eqref{ext:highest-character} applies for
all sufficiently large $\ell$ in each of these finitely many
choices.  Passing to an infinite common subsequence fixes the
cohomological degrees $q_s$ in which the characters occur.

Average a K\"ahler form and the logarithmic weights of the given
line metric over the maximal compact torus $T_c$.  Smooth
semipositivity is preserved.  Hard Lefschetz with the smooth
coefficient $(\ell+1)L_F$~\cite[Theorem~0.1]{DPS2001} gives
$T_c$-equivariant surjections
\begin{equation}\label{ext:lefschetz}
 H^0\bigl(F,\Omega_F^{h-q_s}\otimes(\ell+1)L_F\bigr)
 \longrightarrow H^{q_s}(F,\ell L_F),\qquad h=\dim F.
\end{equation}
The target identification is equivariant because
$K_F+(\ell+1)L_F=\ell L_F$ with the natural linearizations.
Projecting a lift to its character space by weighted Haar averaging
therefore lifts the exact character $\ell w_{P_s}$.

Tensor the lifted nonzero sections with multiplicities $b_s$ and
use the alternating inclusions
$\Omega_F^j\hookrightarrow(\Omega_F^1)^{\otimes j}$ in
characteristic zero.  Their tensor product is nonzero at the generic
point, since $F$ is integral.  By \eqref{ext:cancel} its character
is zero.  Its tensor degree and line exponent are
\[
 p=\sum_s b_s(h-q_s),\qquad
 k=(\ell+1)\sum_s b_s.
\]
The first is fixed and the second tends to infinity.  Invariance
under $T_c$ is equivalent to $T$-invariance in the finite-dimensional
algebraic section representation.  This proves the nontrivial-torus
case.

If $T$ is trivial, Hodge symmetry gives
$\chi(F,\mathcal O_F)=1$.  Riemann--Roch for $\ell L_F$ is a
polynomial with this constant term, so some cohomology group is
nonzero for all sufficiently large $\ell$.  Fix its degree along
an infinite subsequence and apply \eqref{ext:lefschetz}.
For a point take $p=0$ and the constant section for every $k>0$.
\end{proof}

\subsection{Descent before ordinary conversion}

We finish this method by specifying the tensor descent.  Use the
finite-cover structure statement as in
Section~\ref{inv:isometric-descent}: on the universal cover of a
connected finite \'etale cover $\widehat X\to X$ there is an
isometric holomorphic product
\[
 \widetilde X=\mathbb C^a\times R\times F.
\]
Here $R$ is the compact Ricci-flat product, $F$ is projective with
no positive-degree holomorphic forms, and $-K_F$ is smoothly
semipositive.  The chosen deck subgroup acts by translations on
the flat factor, fixes the canonical frame on $R$, and acts on
$F$ through a connected compact torus of holomorphic isometries.
Point factors are allowed.  These are the geometric structure and
frame assertions, independent of the anticanonical nonvanishing
conclusion.

For completeness, the compact action has the algebraic realization
needed by Proposition~\ref{ext:invariant-tensors}.
Since $H^1(F,\mathcal O_F)=0$, it preserves a very ample line up
to isomorphism and hence acts projectively in its complete
linear-system embedding.  The representation is continuous, as
one sees by tracking a projective frame of points on the embedded
variety.  A connected compact abelian subgroup of the projective
linear group is conjugate into a diagonal torus.  Its Zariski
closure is an algebraic torus preserving $F$; quotient by its
ineffective kernel if necessary.  A diagonal projective torus
has a linear lift after fixing one diagonal coordinate, so an
ample linearization is available.  The action on $-K_F$ remains
the intrinsic tangent-determinant action.

Pull each invariant tensor $u_k$ from the proposition to
$\widetilde X$ and use the cotangent inclusion from the product
projection to $F$.  If $\omega_{\mathrm{flat}}$ and $\omega_R$
are the invariant canonical frames of the other factors, then
\[
 u_k\otimes(\omega_{\mathrm{flat}}\otimes\omega_R)^{-k}
 \in H^0\bigl(\widetilde X,
          (\Omega_{\widetilde X}^1)^{\otimes p}
                         \otimes(-kK_{\widetilde X})\bigr)
\]
is nonzero and deck-invariant.  It descends to $\widehat X$ and
is algebraic by projectivity.  Corollary~\ref{conv:tensor-conversion}
now gives a nonzero anticanonical plurisection on this actual
projective finite cover.  Finally the finite \'etale norm recalled
in Section~\ref{inv:isometric-descent} gives a nonzero positive
anticanonical multiple on $X$.  Thus the only invariant objects
needed for this method are the tensors before descent; the
conversion and its Fano-type auxiliary model are ordinary.
